% Options for packages loaded elsewhere
\PassOptionsToPackage{unicode}{hyperref}
\PassOptionsToPackage{hyphens}{url}
%
\documentclass[
  10pt,
]{article}
\usepackage{amsmath,amssymb}
\usepackage{iftex}
\ifPDFTeX
  \usepackage[T1]{fontenc}
  \usepackage[utf8]{inputenc}
  \usepackage{textcomp} % provide euro and other symbols
\else % if luatex or xetex
  \usepackage{unicode-math} % this also loads fontspec
  \defaultfontfeatures{Scale=MatchLowercase}
  \defaultfontfeatures[\rmfamily]{Ligatures=TeX,Scale=1}
\fi
\usepackage{lmodern}
\ifPDFTeX\else
  % xetex/luatex font selection
\fi
% Use upquote if available, for straight quotes in verbatim environments
\IfFileExists{upquote.sty}{\usepackage{upquote}}{}
\IfFileExists{microtype.sty}{% use microtype if available
  \usepackage[]{microtype}
  \UseMicrotypeSet[protrusion]{basicmath} % disable protrusion for tt fonts
}{}
\makeatletter
\@ifundefined{KOMAClassName}{% if non-KOMA class
  \IfFileExists{parskip.sty}{%
    \usepackage{parskip}
  }{% else
    \setlength{\parindent}{0pt}
    \setlength{\parskip}{6pt plus 2pt minus 1pt}}
}{% if KOMA class
  \KOMAoptions{parskip=half}}
\makeatother
\usepackage{xcolor}
\usepackage[margin=22mm]{geometry}
\setlength{\emergencystretch}{3em} % prevent overfull lines
\providecommand{\tightlist}{%
  \setlength{\itemsep}{0pt}\setlength{\parskip}{0pt}}
\setcounter{secnumdepth}{-\maxdimen} % remove section numbering
\ifLuaTeX
\usepackage[bidi=basic]{babel}
\else
\usepackage[bidi=default]{babel}
\fi
\babelprovide[main,import]{italian}
% get rid of language-specific shorthands (see #6817):
\let\LanguageShortHands\languageshorthands
\def\languageshorthands#1{}
\ifLuaTeX
  \usepackage{selnolig}  % disable illegal ligatures
\fi
\usepackage{bookmark}
\IfFileExists{xurl.sty}{\usepackage{xurl}}{} % add URL line breaks if available
\urlstyle{same}
\hypersetup{
  pdflang={it},
  hidelinks,
  pdfcreator={LaTeX via pandoc}}

\author{}
\date{}

\begin{document}

\section{Example project - Facial-recognition access
control}\label{example-project---facial-recognition-access-control}

\begin{quote}
Non-canonical model-validation example.

The example intentionally includes ML and biometric concepts because
they belong to the application domain. The Macro Requirements remain
independent from lower-level architecture choices such as inference
placement, device interconnection and command protocols; those choices
belong to later Decisions.

For comparability with the current ThreatForge project-model convention,
this example uses textual identifiers \texttt{MR-0001} through
\texttt{MR-0004}. The general metamodel still treats the concrete ID
encoding as a deferred Doc-as-Code representation decision.
\end{quote}

\subsection{Project summary}\label{project-summary}

The project controls access to a reserved physical area. A camera
observes a person at the access point, facial recognition is used to
determine whether the person corresponds to an authorized identity, and
access is granted by opening a turnstile when the
project\textquotesingle s access conditions are satisfied.

The MR set describes stable project concerns rather than the progress of
architectural design. Placement of ML inference, network topology,
device protocols and similar solution choices are therefore omitted from
the MR bodies and can be governed later through Decisions and
Requirements.

\subsection{MR-0001 - Accesso controllato all\textquotesingle area
riservata}\label{mr-0001---accesso-controllato-allarea-riservata}

\subsubsection{Intent}\label{intent}

Consentire alle persone autorizzate di accedere a
un\textquotesingle area riservata in modo rapido e controllato, evitando
che il varco si apra automaticamente quando
l\textquotesingle autorizzazione non e stata confermata.

\subsubsection{Context}\label{context}

L\textquotesingle area e protetta da un tornello. Il progetto usa il
riconoscimento del volto come parte della decisione che precede
l\textquotesingle apertura automatica del varco.

\subsubsection{Stakeholders}\label{stakeholders}

\begin{itemize}
\tightlist
\item
  persone autorizzate;
\item
  responsabili dell\textquotesingle area riservata;
\item
  addetti alla sicurezza fisica;
\item
  personale operativo che gestisce il punto di accesso.
\end{itemize}

\subsubsection{Scope - candidate under
test}\label{scope---candidate-under-test}

\begin{itemize}
\tightlist
\item
  il passaggio dal tentativo di ingresso alla decisione di aprire o non
  aprire il varco;
\item
  l\textquotesingle esperienza macro della persona al punto di accesso.
\end{itemize}

\subsubsection{DependsOn}\label{dependson}

\begin{itemize}
\tightlist
\item
  \texttt{MR-0002} Gestione delle persone autorizzate.
\item
  \texttt{MR-0003} Riconoscimento facciale per la verifica della
  persona.
\end{itemize}

\subsection{MR-0002 - Gestione delle persone
autorizzate}\label{mr-0002---gestione-delle-persone-autorizzate}

\subsubsection{Intent}\label{intent-1}

Consentire all\textquotesingle organizzazione di stabilire chi puo
accedere all\textquotesingle area riservata e di mantenere aggiornata
nel tempo l\textquotesingle associazione tra persona, autorizzazione e
informazioni necessarie al riconoscimento.

\subsubsection{Context}\label{context-1}

Le autorizzazioni cambiano nel tempo: una persona puo essere abilitata,
sospesa o non piu autorizzata. Il sistema di riconoscimento deve usare
informazioni coerenti con lo stato corrente deciso
dall\textquotesingle organizzazione.

\subsubsection{Stakeholders}\label{stakeholders-1}

\begin{itemize}
\tightlist
\item
  responsabili dell\textquotesingle area;
\item
  amministratori delle autorizzazioni;
\item
  persone autorizzate;
\item
  personale incaricato dell\textquotesingle onboarding e della revoca.
\end{itemize}

\subsubsection{Scope - candidate under
test}\label{scope---candidate-under-test-1}

\begin{itemize}
\tightlist
\item
  inserimento, aggiornamento, sospensione e revoca delle autorizzazioni;
\item
  associazione tra identita autorizzata e informazioni utilizzate per
  riconoscerla.
\end{itemize}

\subsection{MR-0003 - Riconoscimento facciale per la verifica della
persona}\label{mr-0003---riconoscimento-facciale-per-la-verifica-della-persona}

\subsubsection{Intent}\label{intent-2}

Determinare se il volto osservato al punto di accesso corrisponde a una
persona autorizzata con qualita sufficiente per supportare la decisione
di accesso.

\subsubsection{Context}\label{context-2}

Il progetto utilizza una telecamera e un modello di machine learning per
il riconoscimento facciale. Il modello puo richiedere immagini o
sequenze video adeguate e puo produrre risultati incerti o errati.

\subsubsection{Stakeholders}\label{stakeholders-2}

\begin{itemize}
\tightlist
\item
  persone che si presentano al varco;
\item
  addetti alla sicurezza;
\item
  responsabili del servizio di accesso;
\item
  specialisti che mantengono il sistema di riconoscimento.
\end{itemize}

\subsubsection{Scope - candidate under
test}\label{scope---candidate-under-test-2}

\begin{itemize}
\tightlist
\item
  acquisizione delle informazioni visive necessarie al riconoscimento;
\item
  esecuzione e valutazione del riconoscimento facciale;
\item
  risultato fornito al processo che decide l\textquotesingle accesso.
\end{itemize}

\subsubsection{Assumptions}\label{assumptions}

\begin{itemize}
\tightlist
\item
  il riconoscimento facciale e una capacita prevista dal progetto, non
  una tecnica scelta solo per una metodologia di analisi.
\end{itemize}

\subsection{MR-0004 - Uso responsabile dei dati biometrici e delle
decisioni
automatiche}\label{mr-0004---uso-responsabile-dei-dati-biometrici-e-delle-decisioni-automatiche}

\subsubsection{Intent}\label{intent-3}

Proteggere le persone coinvolte nell\textquotesingle uso del
riconoscimento facciale assicurando che dati biometrici e decisioni
automatiche siano trattati in modo comprensibile, controllabile e
coerente con le finalita di accesso del progetto.

\subsubsection{Context}\label{context-3}

Il riconoscimento facciale utilizza informazioni biometriche e puo
produrre errori che incidono sull\textquotesingle accesso fisico di una
persona. Questi aspetti fanno parte del dominio e del valore del
progetto indipendentemente dalla metodologia di analisi che potra essere
applicata.

\subsubsection{Stakeholders}\label{stakeholders-3}

\begin{itemize}
\tightlist
\item
  persone i cui dati biometrici sono trattati;
\item
  responsabili dell\textquotesingle organizzazione;
\item
  responsabili privacy/compliance applicabili al contesto;
\item
  addetti che gestiscono contestazioni o eccezioni di accesso.
\end{itemize}

\subsubsection{Scope - candidate under
test}\label{scope---candidate-under-test-3}

\begin{itemize}
\tightlist
\item
  trattamento dei dati biometrici utilizzati dal progetto;
\item
  conseguenze e revisione delle decisioni automatiche di riconoscimento
  nel contesto dell\textquotesingle accesso.
\end{itemize}

\subsection{Teachability map}\label{teachability-map}

The MR set is also used to test whether a reader new to the project can
form a simple mental model before seeing Decisions or architecture. A
saved non-canonical Macro Project Map is available in
\texttt{../04-teachability-maps/facial-access-macro-project-map.html}
(with SVG/PNG companions and a textual traceability explanation).

The map distinguishes explicit \texttt{dependsOn} relations from
didactic links inferred only for explanation. These links are not data
flows and do not establish architecture.

\subsection{Why there is no separate MR for PC/LAN/edge
architecture}\label{why-there-is-no-separate-mr-for-pclanedge-architecture}

\texttt{PC}, \texttt{LAN}, inference on the camera, inference on an edge
computer and the protocol used to command the turnstile are solution
choices rather than independent project macro-concerns in this example.
They are candidate \textbf{Decisions} below the relevant MR branches.

A different MR structure would be justified only if a property that
currently looks architectural became an independently valued project
outcome with its own stakeholder meaning. For example, a project-wide
requirement for autonomous operation without external connectivity could
reveal a distinct macro concern. This is a test for the MR boundary, not
temporary status text to place inside an MR.

\subsection{Candidate split test: ML model
lifecycle}\label{candidate-split-test-ml-model-lifecycle}

In this example, model training, deployment, threshold selection and
model update are treated as lower-level concerns of \texttt{MR-0003}
because they serve the single macro intent of facial recognition at the
access point.

A separate MR for \textbf{ML model lifecycle} would be justified only if
the project treats model production/governance as an independent macro
capability with its own stakeholders, macro outcome and Decisions (for
example, if the product itself includes a reusable organizational ML
model-management service).

\subsection{Observation for the future analysis
metamodel}\label{observation-for-the-future-analysis-metamodel}

This example exposes an important distinction that should be preserved
for the later analysis-model study rather than encoded into the MR
structure itself.

Different parts or concerns of one governed project may benefit from
different analysis paradigms:

\begin{itemize}
\tightlist
\item
  \texttt{MR-0003} is a natural candidate for later analysis with an
  AI/ML-oriented threat-analysis method such as STRIDE-AI, because the
  project concern includes an ML-based recognition capability and
  model-specific failure or abuse questions may become relevant;
\item
  \texttt{MR-0004} is a natural candidate for a privacy-oriented
  analysis because it concerns biometric data and effects on people.
  Privacy-specific analysis is outside the thesis implementation scope,
  but the example demonstrates why a general governance model should not
  assume that one analysis methodology covers every relevant concern
  equally well.
\end{itemize}

The working lesson is \textbf{not} that an MR owns one methodology. The
lesson is that later analysis should be able to select a method for an
explicit analysis scope or set of governed subjects. Different methods
may therefore analyze different, overlapping or complementary parts of
the same project.

Their outputs must still converge into one governed
project/documentation lifecycle: findings, documentation gaps, proposed
requirements and other feedback should remain traceable to the same
governed project knowledge and return through common review and
documentation-evolution mechanisms.

This observation is intentionally deferred to the future analysis
metamodel. It must not be used to add STRIDE-AI, privacy or other
method-specific fields to Macro Requirements.

\subsection{What this example tests}\label{what-this-example-tests}

\begin{enumerate}
\def\labelenumi{\arabic{enumi}.}
\tightlist
\item
  Domain-specific technical terms can appear in MR when they are natural
  to the project stakeholders.
\item
  Lower-level architecture choices are omitted from MR bodies rather
  than described through temporary statements such as "not yet decided".
\item
  ML can be represented at MR level without introducing a
  threat-analysis or privacy-analysis methodology.
\item
  Cross-MR dependencies can express collaboration without creating a
  false hierarchy between independent macro concerns.
\item
  A single project can expose concerns suited to different later
  analysis paradigms without partitioning the MR model by methodology.
\item
  Results from different analyses should ultimately return to one
  governed project/documentation lifecycle; the detailed mechanism
  remains a future analysis-metamodel question.
\item
  Intent alone carries macro purpose, value and desired outcome; the
  separate Objective field was removed after duplication testing.
\item
  Scope remains under evaluation and should still be tested for
  redundancy with exclusions/non-goals.
\end{enumerate}

\end{document}
